\SourceLeaf{367}%
lich vorgelegt, und ein kurzer Auszug daraus schon damals in unsern Blättern mitgetheilt (1805 St. 124). Wir bemerken also hier nur, dass dies diejenige Stelle ist, wo Socrates durch ein Beispiel aus der Geometrie anschaulich machen will, wie man sich zur Auflösung einer Aufgabe vorher durch Annahme gewisser näherer Bestimmungen vorzubereiten hat. Die geometrische Aufgabe, welche Socrates hierzu wählt, ist die Frage über die Möglichkeit, ein gegebenes Dreieck in einen gegebenen Kreis einzutragen, aber die Worte, wodurch er erst gewisse Einschränkungen über die Art des Dreiecks festsetzen will, haben den Auslegern viel zu schaffen gemacht. Herr \textsc{Mollweide} führt mit vielem gelehrten Scharfsinn hier aus, dass die dadurch bezeichnete Eigenschaft keine andere ist, als die Zerlegbarkeit des Dreiecks in zwei andere dem Ganzen ähnliche, welches denn freilich im Grunde nichts anders als eine pretiöse Umschreibung des \textit{rechtwinkligen} Dreiecks ist. Die Art wie Hr. M. beweist, dass jene Eigenschaft nur dem rechtwinkligen Dreiecke zukommen kann, ist viel künstlicher und weitläufiger als hier eben nöthig gewesen wäre, da dies gleich unmittelbar aus der Gleichheit der drei Winkel \(ABC,ADB,BDC\) folgt (S. 46).

Die dritte Abhandlung war gleichfalls schon früher unserer Societät handschriftlich vorgelegt, und ein Bericht darüber in unsern gel. Anz. (1807 St. 74) gegeben; sie erscheint hier mit bedeutenden Vermehrungen. Es werden darin zwei von \textsc{Columella} gelehrte Näherungsmethoden erläutert, die Fläche des gleichseitigen Dreiecks und die Fläche eines Kreissegments zu berechnen. Eine kleine Übereilung findet sich S. 71, wo behauptet wird, dass kein anderer Bruch, dessen Zähler und Nenner unter 100 sei, dem wahren Verhältnisse des gleichseitigen Dreiecks zum Quadrate über derselben Seite so nahe kommen könne, als \(\frac{13}{30}\); in der That sind die beiden Brüche \(\frac{29}{67}\) und \(\frac{42}{97}\) genauer.

Der Brief an den verdienten Prof. \textsc{Schneider} in Breslau enthält einige Anmerkungen zu den von \textsc{Hase} in \textsc{Bredows} \textit{Epistolae Parisienses} mitgetheilten Stücken von den freilich sehr unbedeutenden mathematischen Schriften des \textsc{Vitruvius Rufus} und \textsc{Epaphroditus}.
